\section*{Chapter \ref{ch:b3:quantum-angular-momentum} --- Quantum Angular Momentum}

\begin{solution}{exo:b3:quantum-angular-momentum:1}
(a) $[\hat y\hat p_z - \hat z\hat p_y,\ \hat z\hat p_x - \hat x\hat
p_z]$: only the terms sharing a $\hat z, \hat p_z$ pair survive,
giving $\iu\hbar(\hat x\hat p_y - \hat y\hat p_x)$. (b) $[\hat
L^2, \hat L_z] = \sum_i[\hat L_i^2, \hat L_z] = \hat L_x[\hat L_x,
\hat L_z] + [\hat L_x, \hat L_z]\hat L_x + (x \to y)$: the four
terms cancel pairwise. (c) $[\hat L_z, \hat x] = \iu\hbar\hat y$:
$\hat L_z$ generates rotations about $z$ --- of operators as of
states. (d) Common sharpness of two components forces (by the
\omterm{def:b3:quantum-formalism:commutator}{commutator}) sharpness of the third with value zero for all three:
only $l = 0$ achieves it.
\end{solution}

\begin{solution}{exo:b3:quantum-angular-momentum:2}
(a) $\hat L_z = \hbar\operatorname{diag}(1, 0, -1)$. (b) $\hat L_+ =
\hbar\sqrt2\,\begin{pmatrix}0&1&0\\0&0&1\\0&0&0
\end{pmatrix}$, $\hat L_-$ its transpose. (c) $\hat L_x =
\dfrac{\hbar}{\sqrt2}\begin{pmatrix}0&1&0\\1&0&1\\0&1&0
\end{pmatrix}$: characteristic polynomial $\lambda(\lambda^2 -
\hbar^2)$. (d) The $L_x = +\hbar$ eigenvector is $\tfrac12(1,
\sqrt2, 1)$: probabilities $\tfrac14, \tfrac12, \tfrac14$ for $m =
+1, 0, -1$.
\end{solution}

\begin{solution}{exo:b3:quantum-angular-momentum:3}
(a) $I = \mu r_0^2 = \qty{1.46e-46}{kg.m^2}$; $B = \hbar^2/2I =
\qty{3.8e-23}{J} = \qty{0.24}{meV}$. (b) $2B/h = \qty{115}{GHz}$,
$\lambda = \qty{2.6}{mm}$. (c) \qty{231}{GHz} and \qty{346}{GHz}.
(d) The spacing gives $B$, hence $I = \mu r_0^2$: the bond length of
a molecule light-years away, to roughly half the spacing's relative
precision.
\end{solution}

\begin{solution}{exo:b3:quantum-angular-momentum:4}
(a) $m\hbar$ with $m = -2 \dots 2$. (b) $\cos\theta = 2/\sqrt6$:
$\theta = 35.3^\circ$. (c) Perfect alignment would make $L_x = L_y
= 0$ sharp together with $L_z$ --- forbidden by the \omterm{def:b3:quantum-formalism:commutator}{commutators}
except in the trivial $l = 0$ case. (d) $l/\sqrt{l(l+1)} \to 1$: for
large $l$ the cone closes onto the axis and the classical arrow
returns.
\end{solution}

\begin{solution}{exo:b3:quantum-angular-momentum:5}
(a) With no $\varphi$ dependence, the operator gives
$-\hbar^2(\sin\theta)^{-1}\partial_\theta(\sin\theta\,(-\sin\theta))
= 2\hbar^2\cos\theta$. (b) The same computation with the
$\eu^{\pm\iu\varphi}$ factor: \omterm{def:b3:quantum-formalism:observable}{eigenvalue} $2\hbar^2$; $\hat L_z$
gives $\pm\hbar$. (c) $\cos\theta$ and $\sin\theta\,\eu^{\pm\iu
\varphi}$ change sign under $\vect r \to -\vect r$ ($\theta \to \pi
- \theta$, $\varphi \to \varphi + \pi$): parity $-1 = (-1)^1$. (d)
They are eigenvectors of the Hermitian $\hat L_z$ with different
\omterm{def:b3:quantum-formalism:observable}{eigenvalues}.
\end{solution}

\begin{solution}{exo:b3:quantum-angular-momentum:6}
(a) $2J + 1$ states share the level (degeneracy); the Boltzmann
factor taxes its energy. (b) Maximise the product: $\dd/\dd J = 0$
gives $2 = (2J+1)^2B/k_{\text{B}}T$, the stated $J_{\max}$. (c) At
\qty{10}{K}: $J_{\max} \approx 0.8$, so $J = 1$ dominates and the
\qty{115}{GHz} and \qty{230}{GHz} lines shine; at \qty{300}{K}:
$J_{\max} \approx 7$. (d) $T \approx 2B(J_{\max} + \tfrac12)^2/
k_{\text{B}} \approx \qty{310}{K}$ --- a rotational thermometer.
\end{solution}

\begin{solution}{exo:b3:quantum-angular-momentum:7}
(a) The upper level splits into five, the lower into three, all
with the same spacing $\mu_{\text{B}}B$. (b) With equal spacings,
$h\nu = h\nu_0 + (\Delta m)\mu_{\text{B}}B$ and $\Delta m \in \{0,
\pm1\}$: three positions only. (c) $\mu_{\text{B}}/h = \qty{14}{GHz/T}$, so at $B = \qty{2}{T}$
the splitting is \qty{28}{GHz}, i.e.\ $\Delta\lambda =
\lambda^2\Delta\nu/c \approx \qty{23}{pm}$ --- resolvable since
Zeeman's 1896 gratings.
(d) The electron's own spin, with its anomalous factor $g \approx
2$: the ``anomalous'' patterns were spin's fingerprints before spin
was named (\cref{ch:b3:spin-two-level}).
\end{solution}

\begin{solution}{exo:b3:quantum-angular-momentum:8}
In the stated units $U_{\text{eff}} = -2/x + l(l+1)/x^2$. (a) For
$l = 1$: minimum at $x = 2$ ($r = 2a_0$), depth $-E_{\text{I}}/2$.
(b) $U_{\text{eff}} > 0$ for $x < 1$: inside one Bohr radius the
wall wins. (c) $\psi \sim r^l$ near the origin: only $l = 0$ states
have nonzero density at $r = 0$. (d) Capturing an electron requires
electron density \emph{at} the nucleus: the $s$ electrons, alone
unbarred by the centrifugal wall, are the ones swallowed.
\end{solution}

\begin{solution}{exo:b3:quantum-angular-momentum:9}
(a) $\hat L_x = (\hat L_+ + \hat L_-)/2$ changes $m$: diagonal
elements vanish. (b) By symmetry the two transverse squares are
equal, and their sum is $\langle\hat L^2 - \hat L_z^2\rangle$. (c)
At $m = l$: $\Delta L_x\Delta L_y = l\hbar^2/2 = \tfrac\hbar2
\,l\hbar$: equality --- the stretched state is as aligned as the
algebra permits. (d) The cone: sharp $L_z$, vanishing transverse
means, equal transverse spreads --- a vector of definite length and
projection, democratically smeared in azimuth.
\end{solution}

\begin{solution}{exo:b3:quantum-angular-momentum:10}
(a) Energy bookkeeping with $E_J = BJ(J+1)$: $J \to J + 1$ adds
$2B(J+1)$; $J \to J - 1$ subtracts $2BJ$. (b) $\Delta J = 0$ is
forbidden, and the two branches start at $\pm 2B$: nothing falls at
the pure vibrational frequency. (c) Two combs of spacing $2B$
flanking a gap, intensities rising to $J_{\max}$ then falling. (d)
The band's width is the populated extent of the combs, growing as
$\sqrt{T}$: those thermally populated wings are precisely where a
saturated greenhouse band keeps absorbing.
\end{solution}

\begin{solution}{exo:b3:quantum-angular-momentum:11}
(a) $\|\hat J_\pm\ket{j,m}\|^2 = \bra{j,m}\hat J_\mp\hat J_\pm
\ket{j,m} = \hbar^2[j(j+1) - m(m\pm1)]$. (b) Climbing past $m = j$
would give $j(j+1) - j(j+1) = 0$ then negative norms one step
further: the chain must contain the annihilated top state. (c) The
top and bottom differ by an integer number of unit steps: $2j \in
\mathbb N$: $j = 0, \tfrac12, 1, \tfrac32, \dots$ (d) The algebra
never uses wave functions; only the demand that
$\eu^{\iu m\varphi}$ be single-valued on the circle forces integer
$m$ --- a condition binding orbital motion and leaving intrinsic
(spin) angular momenta free to be half-integer.
\end{solution}

\begin{solution}{exo:b3:quantum-angular-momentum:12}
(a) $n\hbar$ per unit volume. (b) $\omega = n\hbar V/(\tfrac12 Ma^2)
= 2n\hbar/\rho a^2 = 2 \times \num{8.5e28} \times \num{1.055e-34}/
(7900 \times \num{e-4}) \approx \qty{2.3e-5}{rad/s}$ --- invisible
in one shot. (c) Reversing in step with the fibre's resonance
accumulates $\sim Q$ kicks: $\omega \sim \qty{2e-3}{rad/s}$,
milliradian swings on a period of seconds --- visible with a mirror
and a light beam, as Einstein and de Haas saw. (d) The measured
gyromagnetic ratio was $e/m_{\text{e}}$, twice the orbital
$e/2m_{\text{e}}$: iron's magnetism is not orbital currents but
electron \emph{spin}, whose $g \approx 2$ the next chapters explain.
\end{solution}

\begin{solution}{pb:b3:quantum-angular-momentum:1}
\textbf{1.} $E_J = BJ(J+1)$, degeneracy $2J + 1$.
\textbf{2.} $I = \mu r_0^2 = \qty{1.46e-46}{kg.m^2}$: $B =
\hbar^2/2I = \qty{3.8e-23}{J}$, i.e.\ $B/h = \qty{57.7}{GHz}$ ---
the measured constant.
\textbf{3.} $115$, $231$, \qty{346}{GHz}.
\textbf{4.} $E_{J+1} - E_J = 2B(J+1)$ grows linearly: consecutive
lines differ by the constant $2B$. The spacing gives $I$, hence the
bond length --- structural chemistry by radio.
\textbf{5.} $\lambda = c/\nu = \qty{2.6}{mm}$. Atmospheric water
vapour absorbs millimetre waves greedily; hence Atacama, Mauna Kea,
the South Pole --- high, cold, dry.
\textbf{6.} The photon is a spin-$1$ particle carrying $\hbar$:
total angular momentum conservation obliges the molecule's $J$ to
change by one unit.
\textbf{7.} Its charge distribution is symmetric at every rotation
angle: no oscillating dipole, no dipole line --- perfect
camouflage.
\textbf{8.} $B_{\text{H}_2} = \hbar^2/2\mu r_0^2 \approx
\qty{7.6}{meV}$; the symmetric molecule's weak quadrupole
transitions require $\Delta J = 2$, costing $6B \approx
\qty{46}{meV}$ --- fifty times $k_{\text{B}}T$ at \qty{10}{K}: cold
hydrogen cannot radiate at all.
\textbf{9.} $2B = \qty{0.48}{meV} < k_{\text{B}}T = \qty{0.86}{meV}$:
collisions keep the first rungs populated --- the ladder works
exactly where the clouds live.
\textbf{10.} $n_1/n_0 = 3\,\eu^{-0.478/0.862} = 1.7$: the emitting
level is well stocked.
\textbf{11.} $J_{\max} \approx \sqrt{0.86/0.48} - 0.5 \approx 0.8$:
the population peaks at $J = 1$.
\textbf{12.} CO delivers the photons; converting them to total gas
mass rests on an assumed CO-to-H$_2$ abundance --- luminous
messenger, calibrated ransom.
\textbf{13.} $\Delta\nu = \qty{115}{MHz}$: $v = c\,\Delta\nu/\nu
\approx \qty{300}{km/s}$, receding (frequency lowered) --- galactic
orbital speeds.
\textbf{14.} $\pm\qty{2}{km/s}$ spans \qty{1.5}{MHz}; thermal
motion at \qty{10}{K} gives only $\sim\qty{40}{kHz}$: the width is
turbulence, not temperature --- linewidths map a cloud's internal
weather.
\textbf{15.} Around a central mass $v \propto 1/\sqrt r$ should
fall (Kepler); the measured flatness means the enclosed mass keeps
growing with radius --- invisible matter enveloping the luminous
disc: dark matter.
\textbf{16.} The ratio of populations, i.e.\ the excitation
temperature of the gas.
\textbf{17.} At \qty{57.6}{GHz}: cosmic expansion has doubled every
wavelength in flight --- the cosmological redshift the last chapter
of this book returns to.
\textbf{18.} Photon flux $\to$ CO luminosity (distance needed)
$\to$ CO count (excitation model) $\to$ H$_2$ mass (the
CO-to-H$_2$ ``X-factor'': the weakest link, calibrated locally and
exported cautiously).
\textbf{19.} A \qty{10}{K} core emits nothing an optical telescope
can see; its rotational lines are simultaneously its only
thermometer, speedometer and scale.
\textbf{20.} A thermometer is sensitive where level spacing $\sim
k_{\text{B}}T$: \qty{5.5}{K} spacing responds strongly across
$5$--$50$ K, while a \qty{2}{eV} optical level would be populated
by nothing at all.
\textbf{21.} The line saturates (optically thick) and, worse, CO
freezes out onto dust grains in the densest, coldest gas:
astronomers switch to rarer isotopologues ($^{13}$CO, C$^{18}$O)
and other molecules.
\textbf{22.} Each species needs its own density and temperature to
shine: together they tomograph the cloud --- outer envelopes in CO,
dense cores in NH$_3$ and HCN.
\textbf{23.} A line is seen in contrast against the cosmic
background; as the gas temperature approaches \qty{2.7}{K} its
excitation equilibrates with the background and the contrast
vanishes: no colder cloud can be read this way.
\textbf{24.} $\lambda/20 = \qty{130}{\micro m}$ against
$\qty{25}{nm}$ for optical work: five thousand times more
forgiving --- which is why twelve-metre millimetre dishes stand in
the open desert while optical mirrors live in domes.
\textbf{25.} A ladder $BJ(J+1)$ with $B/h = \qty{57.6}{GHz}$, alive
at ten kelvin and read at \qty{2.6}{mm}, measures the mass,
temperature, turbulence and rotation of the cold universe --- and
its flat rotation curves are a standing exhibit for dark matter.
\end{solution}
